\section{Introduction}
\label{sec:intro}

Assessing the security posture of large computing infrastructures is
notoriously difficult because they are complex systems, dynamic in their
connectivity and usage patterns, and they often carry out processes that
should not be disrupted. Performing security analysis directly on a
production system can cause problems: a single aggressive probe against a
database replica, a message broker, or an authentication service can cascade
into an outage that violates a regulated availability requirement, corrupts
state, or interrupts a process on which people depend.

Because of these risks, the analysis of the security posture of an
infrastructure is usually performed manually and in an ad hoc fashion, by
analysts whose expertise is scarce and expensive. As a result, these
assessments are carried out sporadically, and the evolution of the network
topology, the introduction of a new service, or a change in the configuration
of a component might introduce a vulnerability that invalidates the results
of the previous analysis. In the meantime, adversaries that are increasingly
augmented with AI operate continuously, automatically, and at
scale~\cite{measuring}. The proposed research targets this widening gap.

\textbf{The goal of the proposed research is to develop a new agentic
framework, called \sysname, to support the continuous assessment of the
security posture of computing infrastructures based on open-source
software.} \sysname uses cooperating teams of AI agents to capture the
configuration of a target system, monitor its execution, create a
\emph{digital twin} of the system, run red team exercises against the twin,
and guide a human-driven verification and remediation of the issues that are
found in the original system. The core principle of \sysname is a strict
separation between \emph{observation}, which is passive, non-disruptive, and
the only activity that ever touches the production system, and
\emph{analysis}, which is aggressive, potentially destructive, and confined
entirely to the digital twin. This approach is feasible because the components of the infrastructure are
open source, and therefore their behavior, configuration surface, and version
history are transparent and inspectable. This transparency makes it possible to run a
functionally equivalent instance of a component in a controlled environment,
something generally infeasible for closed, proprietary appliances.

\paragraph{A motivating scenario.}
Consider a regional health-care provider whose electronic health-record
system runs on a stack of open-source components: a
PostgreSQL~\cite{postgresql} cluster with streaming replication, a set of
application servers behind an nginx~\cite{nginx} reverse proxy, a
Keycloak~\cite{keycloak} identity provider federated with the hospital
directory, a RabbitMQ~\cite{rabbitmq} message broker connecting clinical
subsystems, and a fleet of medical-device gateways that speak proprietary
protocols. This deployment is subject to strict availability and privacy
requirements, and the security team cannot fuzz the application
servers, cannot flood the broker with malformed messages, and cannot run an
exploit against the identity provider. As a result, the questions that matter to the security team remain
unanswered. Can an attacker who compromises a
single device gateway pivot through the broker to reach patient records? Does
the proxy configuration expose an internal administrative endpoint through a
path that the access rules did not anticipate? Each of these questions concerns the \emph{deployment} as a whole, and none can be answered from an advisory or a source repository. \sysname is designed for this situation: it reconstructs the deployment as a
digital twin, exercises exactly these attack paths against the twin, and
reports the confirmed chains to the security team for verification on the
production system, all without ever sending a single aggressive packet to the
live infrastructure.

A possible alternative to a digital twin would be to leverage the staging
environment that engineers already use to test applications before
deployment. In practice, this would not be sufficient, as staging
environments are built and maintained manually, they drift from production as
the infrastructure evolves, and they are typically provisioned for functional
testing rather than to preserve the security-relevant structure of the
deployment: the exact component versions, the network reachability relations,
the trust and authentication dependencies, and the data that determines
authorization decisions. A staging environment that omits the device gateways, or that seeds an empty user directory, will not exhibit the attack paths exploitable in the real infrastructure. \sysname constructs a digital twin whose fidelity is measured and controlled
with respect to explicit security properties, as described in
Section~\ref{sec:thrust2}, and that is kept synchronized with the production
system by a continuous capture-and-observe loop, as described in
Section~\ref{sec:thrust1}.

\paragraph{Alignment with PESOSE Track~3 and contributions.}
This project is submitted to NSF~26-506, Track~3, which addresses
vulnerabilities in open-source products and infrastructure, including both
technical risks and socio-technical ones. The ecosystem we target is anchored on two existing, widely deployed open-source products, described in Section~\ref{sec:ecosystem}. The first is \textbf{nginx}~\cite{nginx,nginxsrc}, which we will harden against the memory-safety, configuration, and request-smuggling defect classes catalogued in Section~\ref{sec:landscape}. The second is \textbf{ZMap}~\cite{zmap,zmapsrc}, which we will extend into an instrument for safe, budgeted, and information-rich discovery inside production networks. Every finding that we confirm is driven upstream, so the work improves both
individual deployments and the shared ecosystem on which they are built.

The project contributes: (i) agentic workflows that capture the configuration
of an infrastructure into a continuously maintained knowledge base, together
with upstream extensions that make open-source discovery safe on production
networks; (ii) techniques for replicating large infrastructures at reduced
scale, supported by a \emph{quantitative fidelity measure}; (iii) methods for
chaining vulnerabilities across a networked infrastructure, extending our
prior work on AI-planning-based privilege escalation and on a universal
theory of in-program exploitation to the setting of complete ecosystems; and
(iv) the hardening of nginx and ZMap, a public corpus of reference
infrastructures with ground truth, and an open-source release of the
framework.
